%!TEX root = ../chapter06.tex 
%----------------------------------------------------------------------------------------
%	
%---------------------------------------------------------------------------------------- 
%----------------------------------------------------------------------------------------
%	EXERCICES
%----------------------------------------------------------------------------------------

\newpage 
\loadgeometry{widelayout}  
\pagestyle{scrheadings} % Print headers again  
\KOMAoptions{
	headwidth=textwithmarginpar,
	headsepline=true,
	headsepline= 0.5pt,
	footwidth=textwithmarginpar,
} 
\onehalfspacing
\subsection{Exercices}
\begin{exercise} Dans le repère oblique $(O;I,J)$ ci-dessous, préciser les coordonnées des points suivants :
	
	\noindent\parbox{0.3\linewidth}{ 
		$O(\qquad;\qquad )$\\
		$I(\qquad;\qquad )$\\
		$J(\qquad;\qquad )$\\ 
		$A( \qquad;\qquad )$\\  
		$B(\qquad;\qquad)$\\  
		$C(\qquad;\qquad)$\\  
		$D(\qquad;\qquad)$\\ 
		$E(\qquad;\qquad)$
	}  \hfill \parbox{0.7\linewidth}{
		\begin{tikzpicture}[scale=0.9] 	
			
			\tkzSetUpPoint[size=3, fill=black]
			\tkzfctset{tan style/.style={->,>=latex, color=red, line width=1.2pt}}   
			\tkzSetUpLine[line width= 1.2pt, add=.5 and .5] 
			
			\def \xmin{-4}; 
			\def \xmax{7.5}; 
			\def \ymin{-3};
			\def \ymax{5};	
			\def \alpha{-15}
			\def \beta{55}
			\def \xnorm{1.2};
			\def \ynorm{0.8}; 
			\pgftransformcm{\xnorm*cos(\alpha)}{\xnorm*sin(\alpha)}{\ynorm*cos(\beta)}{\ynorm*sin(\beta)}{\pgfpointorigin}
			\path[clip,preaction = {draw=gray}] 
			({(\xmax*\ynorm*sin(\beta)-\ymax*\ynorm*cos(\beta))/(\xnorm*\ynorm*sin(\beta-\alpha))},{(-\xmax*\xnorm*sin(\alpha)+\ymax*\xnorm*cos(\alpha))/(\xnorm*\ynorm*sin(\beta-\alpha))}	) -- ( 
			{(\xmax*\ynorm*sin(\beta)-\ymin*\ynorm*cos(\beta))/(\xnorm*\ynorm*sin(\beta-\alpha))},{(-\xmax*\xnorm*sin(\alpha)+\ymin*\xnorm*cos(\alpha))/(\xnorm*\ynorm*sin(\beta-\alpha))}) -- (
			{(\xmin*\ynorm*sin(\beta)-\ymin*\ynorm*cos(\beta))/(\xnorm*\ynorm*sin(\beta-\alpha))},{(-\xmin*\xnorm*sin(\alpha)+\ymin*\xnorm*cos(\alpha))/(\xnorm*\ynorm*sin(\beta-\alpha))}) -- (
			{(\xmin*\ynorm*sin(\beta)-\ymax*\ynorm*cos(\beta))/(\xnorm*\ynorm*sin(\beta-\alpha))},{(-\xmin*\xnorm*sin(\alpha)+\ymax*\xnorm*cos(\alpha))/(\xnorm*\ynorm*sin(\beta-\alpha))}) --   cycle; 
			\tkzInit[xmin=-8, ymin=-8,  xmax=8.5,ymax=9.5]
			\begin{scope}[dashed] 
				\tkzGrid[ xstep=1, ystep=1](-11,-9)(10,15)
			\end{scope}  
			\tkzDrawX[left=5pt, label=, font=\scriptsize, line width=1.2pt ]  
			\tkzDrawY[below =5pt, label=, font=\scriptsize, line width=1.2pt] 
			\tkzRep[color  = Maroon,xlabel=,ylabel=, font=\scriptsize, line width = 2pt,xnorm=1, ynorm=2]   
			
			\tkzDefPoint(0,0){O} 
			\tkzDefPoint(1,0){I} 
			\tkzDefPoint(0,2){J} 
			\tkzDefPoint(4,5){A} 
			\tkzDefPoint(4,0){B} 
			\tkzDefPoint(0,3){C} 
			\tkzDefPoint(-3,4){D}
			\tkzDefPoint(-1,-3){E}  
			\tkzDrawPoints(O,I,J,A,B,C,D,E)
			\tkzLabelPoints[below,  font=\scriptsize](O,I,J)
			\tkzLabelPoints[below, font=\scriptsize](A,B,C,D,E) 
		\end{tikzpicture}
	} 
\end{exercise}
\begin{eBox}
	Dans les exercices suivants, le repère $(O;I,J)$ est orthonormé.
\end{eBox} 
\begin{exercise}  
	Soit les points $A(3;4)$, $B(3;-2)$ et $C(-5;-2)$.
	\begin{enumerate}[leftmargin=*,label=\alph*)]
		\item Préciser un segment horizontal, un segment vertical et donner leurs longueurs.
		\item Calculer la longueur du segment restant.
		\item Calculer les angles du triangle au dixième de degré près.
	\end{enumerate}
\end{exercise} 
\begin{exercise}  
	Soit les points $A(1;8)$, $B(4;1)$ et $C(7;1)$ dans un repère orthonormé $(O;I,J)$.
	
	Calculer le périmètre  et l'aire du triangle $ABC$. 
\end{exercise}    
\begin{exercise}\label{chapter06:exercice:distances:naturetriangle01}
	Soit $A(4;-11)$, $B(10;-10)$ et $C(2;1)$.
	\begin{enumerate}[label=\alph*),leftmargin=*]
		\item Calculer $AB^2$, $AC^2$ et $BC^2$
		\item Écrire les longueurs  $AB$, $AC$ et $BC$ sous la forme $a\sqrt{k}$, $a,k\in\mathbb{N}$, $k$ le plus petit possible.
		\item Préciser et justifier la nature du triangle $ABC$ (équilatéral non isocèle, isocèle, rectangle).
	\end{enumerate}
\end{exercise}
\begin{exercise}\label{chapter06:exercice:distances:naturetriangle02}
	Même questions pour les points $A(1;3)$, $B(13;19)$ et $C(17;15)$. 
	%	Soit $A(-1;2)$, $B(-1;2)$ et $C(3;0)$.  
\end{exercise} 
\begin{exercise}\label{chapter06:exercice:distances:naturetriangle03}
	Même questions pour les points $A(1;20)$, $B(21;35)$ et $C(-14;40)$.  
\end{exercise}  
%\begin{exercise}
%Soit $\mathscr{C}$ le cercle de centre $A$ et de diamètre $[BC]$. Donner la valeur du rayon pour $A(-2;9)$ et $B(-\numprint{1.5};5)$.
%\end{exercise}

\begin{exercise}\label{chapter06:exercice:distances:cercle01} Déterminer le rayon du cercle de diamètre $[BC]$ pour $B(-6;-1)$ et $C(-2;7)$.
\end{exercise}
\begin{exercise}\label{chapter06:exercice:distances:cercle02}  
	Déterminer si le point $B(-6;5)$ est intérieur, extérieur ou appartient au cercle $\mathscr{C}$ de centre $A(8;-2)$ de rayon $\tfrac{7}{3}\sqrt{15}$.
\end{exercise}
\begin{exercise}  
	Soit $A(12;6)$ et $M(-12;24)$.  
	Montrer que la droite $(AM)$ est tangente au cercle de centre $O$ passant par $A$, et préciser la distance entre $O$ et $(AM)$
\end{exercise} 
\begin{exercise} \label{chapter06:exercice:distances:entrainement}  	Calculer $PQ$  dans chacun des cas suivants. Exprimer la longueur sous la forme d'un entier ou $a\sqrt{k}$, avec $k$ un entier le plus petit possible, et $a$ entier.
	\begin{multicols}{3}
		\begin{enumerate}[label=\alph*),leftmargin=* ]
			\item $P(2;1)$ et $Q(1;5)$
			\item $P(-2;1)$ et $Q(3;7)$
			\item $P(-3;-4)$ et $Q(-6;-7)$
			\item $P(2;7)$ et $Q(-1;-2)$ 
			\item $P(9;5\sqrt{3})$ et $Q(2;7\sqrt{3})$ 
			\item $P(7\sqrt{2};4\sqrt{2})$ et $Q(3\sqrt{2};-4\sqrt{2})$ 
		\end{enumerate}
	\end{multicols}
\end{exercise} 
%----------------------------------------------------------------------------------------
%
%----------------------------------------------------------------------------------------
\newpage 

\begin{proof}[solution de l'exercice \ref{chapter06:exercice:distances:naturetriangle01}]   
	$AB=\sqrt{37}$. $AC=2\sqrt{37}$. $BC=\sqrt{185}$. Triangle   rectangle. 
\end{proof}
\begin{proof}[solution de l'exercice \ref{chapter06:exercice:distances:naturetriangle02}] \hfill 
	$AB=AC=20$. $BC=4\sqrt{2}$. Triangle   isocèle. 
\end{proof}
\begin{proof}[solution de l'exercice \ref{chapter06:exercice:distances:naturetriangle03}]   
	$AB=AC=25$. $BC=25\sqrt{2}$. Triangle rectangle isocèle. 
\end{proof}
\begin{proof}[solution de l'exercice \ref{chapter06:exercice:distances:cercle01} ]   
	$r=2\sqrt{5}$.
\end{proof}
\begin{proof}[solution de l'exercice \ref{chapter06:exercice:distances:cercle02}]   
	$BC=7\sqrt{5}$. Point extérieur au cercle.
\end{proof}
\begin{proof}[solution de l'exercice \ref{chapter06:exercice:distances:entrainement}] \hfill 
	\begin{multicols}{3}
		\begin{enumerate}[leftmargin=*,label=\alph*)]
			\item $PQ=\sqrt{17}$
			\item $PQ=\sqrt{61}$
			\item $PQ=2\sqrt{3}$
			\item $PQ=10$
			\item $PQ=\sqrt{61}$
			\item $PQ=4\sqrt{10}$
		\end{enumerate}
	\end{multicols}
\end{proof}

